\section*{第 \ref{ch:b2:integration} 章——积分}

\begin{solution}{exo:b2:integration:1}
$\int_0^1 \frac{\dd t}{\sqrt{t(1-t)}}$：在 $0$ 附近，$\sim t^{-1/2}$
（$\alpha = \frac12 < 1$：收敛）；在 $1$ 附近，$\sim
(1-t)^{-1/2}$：收敛。故积分收敛（其值为 $\pi$，由
换元 $t = \sin^2\theta$ 得）。

$\int_1^\infty \frac{\ln t}{t^2}$：$\frac{\ln t}{t^2} = o(t^{-3/2})$：
收敛（分部积分可得其值为 $1$）。

$\int_0^\infty \frac{\dd t}{1 + t^2\sin^2 t}$：发散。在 $t =
n\pi$ 附近，记 $t = n\pi + u$，则 $\sin^2 t = \sin^2 u \leq u^2$，于是在
$\abs u \leq \frac{1}{n}$ 上，$1 + t^2\sin^2 t \leq 1 + (n\pi +
1)^2u^2 \leq C n^2 u^2 + 1$；因此
\[
\int_{n\pi - 1/n}^{n\pi + 1/n} \frac{\dd t}{1 + t^2\sin^2 t}
\geq \int_{-1/n}^{1/n} \frac{\dd u}{1 + Cn^2u^2}
= \frac{2\arctan\sqrt C}{\sqrt C}\cdot\frac{1}{n} ,
\]
这是一个发散的调和型级数的项，对 $n$ 求和时
原函数无界。
\end{solution}

\begin{solution}{exo:b2:integration:2}
作换元 $u = \lambda t$：
\[
\int_0^\infty t^n \eu^{-\lambda t}\dd t
= \frac{1}{\lambda^{n+1}}\int_0^\infty u^n\eu^{-u}\dd u
= \frac{\Gamma(n+1)}{\lambda^{n+1}} = \frac{n!}{\lambda^{n+1}} .
\]
取 $t = \eu^{-u}$（$\dd t = -\eu^{-u}\dd u$）：
\[
\int_0^1 (\ln t)^n \dd t = \int_0^{\infty} (-u)^n \eu^{-u}\,\dd u
= (-1)^n\, n! .
\]
\end{solution}

\begin{solution}{exo:b2:integration:3}
收敛性：被积函数 $\leq \frac{1}{1+t^2}$（在 $\infty$ 附近），
在 $0$ 附近有界（两个因子在远离 $0$ 时有正下界）：
对每个 $x$ \omterm{def:b2:series:def}{绝对收敛}。作换元 $t = \frac1u$
（$\dd t = -\frac{\dd u}{u^2}$）：
\[
I(x) = \int_0^\infty \frac{1}{\bigl(1 +
\frac1{u^2}\bigr)\bigl(1 + u^{-x}\bigr)}\cdot\frac{\dd u}{u^2}
= \int_0^\infty \frac{u^x}{(1 + u^2)(1 + u^x)}\,\dd u .
\]
将 $I(x)$ 的两种表达式相加得：
\[
2I(x) = \int_0^\infty \frac{1 + t^x}{(1+t^2)(1+t^x)}\dd t
= \int_0^\infty \frac{\dd t}{1 + t^2} = \frac{\pi}{2} :
\]
故 $I(x) = \frac\pi4$，与 $x$ 无关。
\end{solution}

\begin{solution}{exo:b2:integration:4}
在 $0^+$ 附近，记 $u = \abs{\ln t} \to \infty$。若 $\alpha < 1$，
则不论 $\beta$ 取何值均收敛（与 $t^{-\alpha'}$ 比较，
其中 $\alpha < \alpha' < 1$，对数因子被压制）。若 $\alpha > 1$，
则不论 $\beta$ 取何值均发散（与 $t^{-\alpha''}$ 比较，
其中 $1 < \alpha'' < \alpha$）。若 $\alpha = 1$，作换元 $t =
\eu^{-u}$：
\[
\int_0^{1/2} \frac{\dd t}{t\,\abs{\ln t}^\beta}
= \int_{\ln 2}^{\infty} \frac{\dd u}{u^\beta},
\]
当且仅当 $\beta > 1$ 时收敛。综上，收敛当且仅当 $\alpha < 1$，
或（$\alpha = 1$ 且 $\beta > 1$）——恰与贝特朗
级数形成对照。
\end{solution}

\begin{solution}{exo:b2:integration:5}
在 $\intco{0}{\infty}$ 上的\omterm{def:b2:metric:continuity}{连续}性：以
$\bigl|\frac{\eu^{-xt}}{1+t^2}\bigr| \leq \frac{1}{1+t^2}$ 为控制函数，
它可积，且关于 $x \geq 0$ 一致成立，
由\cref{thm:b2:integration:continuity}可得。

$C^2$ 在 $\intoo{0}{\infty}$ 上：当 $x \geq a > 0$ 时，关于 $x$ 的前两阶
导数 $\frac{-t\,\eu^{-xt}}{1+t^2}$ 与
$\frac{t^2\eu^{-xt}}{1+t^2}$ 分别被 $t\,\eu^{-at}$ 与
$\eu^{-at}$ 控制，\omterm{def:b2:metric:continuity}{连续}两次应用
\cref{thm:b2:integration:leibnizrule}。于是
\[
F''(x) + F(x) = \int_0^\infty \frac{t^2 + 1}{1 + t^2}\,\eu^{-xt}\dd
t = \int_0^\infty \eu^{-xt}\dd t = \frac1x .
\]
极限：$0 \leq F(x) \leq \int_0^\infty \eu^{-xt}\dd t = \frac1x \to
0$。
\end{solution}

\begin{solution}{exo:b2:integration:6}
在 $\intcc{\varepsilon}{M}$ 上，
于两段中分别作换元 $u = at$ 与 $u = bt$，即：
\[
\int_\varepsilon^M \frac{f(at) - f(bt)}{t}\dd t
= \int_{a\varepsilon}^{aM}\frac{f(u)}{u}\dd u
- \int_{b\varepsilon}^{bM}\frac{f(u)}{u}\dd u
= \int_{a\varepsilon}^{b\varepsilon} \frac{f(u)}{u}\dd u
- \int_{aM}^{bM} \frac{f(u)}{u}\dd u .
\]
第一项：$f(u) = f(0) + o(1)$（$0$ 附近），且 $\int_{a\varepsilon}
^{b\varepsilon} \frac{\dd u}{u} = \ln\frac ba$，该项趋于
$f(0)\ln\frac ba$。第二项：$f(u) \to f(\infty)$，计算相同，
趋于 $f(\infty)\ln\frac ba$。因此该\omterm{def:b2:integration:improper}{广义积分}
收敛于 $\bigl(f(0) - f(\infty)\bigr)\ln\frac ba$。

取 $f(t) = \eu^{-t}$（$f(0) = 1$，$f(\infty) = 0$），$a = 1$，$b
= 2$：
\[
\int_0^\infty \frac{\eu^{-t} - \eu^{-2t}}{t}\dd t = \ln 2 .
\]
\end{solution}

\begin{solution}{exo:b2:integration:7}
将被积函数延拓为 $0$（即 $t = n$ 以外的部分）：$g_n(t) = (1 -
\frac tn)^n t^{x-1}\mathbf{1}_{t \leq n}$。逐点地，$g_n(t) \to
\eu^{-t}t^{x-1}$（复利极限，见本科一年级分册）。
控制：由 $\ln(1 - u) \leq -u$ 得 $(1 - \frac tn)^n \leq
\eu^{-t}$ 在 $\intcc{0}{n}$ 上成立，于是 $\abs{g_n(t)} \leq
\eu^{-t}t^{x-1} = \varphi(t)$，可积。控制收敛定理：
\[
\int_0^n \Bigl(1 - \frac tn\Bigr)^n t^{x-1}\dd t
\xrightarrow[n\to\infty]{} \int_0^\infty \eu^{-t}t^{x-1}\dd t
= \Gamma(x) .
\]
（对左端反复分部积分，可得欧拉乘积
形式 $\Gamma(x) = \lim \frac{n!\,n^x}{x(x+1)\cdots(x+n)}$。）
\end{solution}

\begin{solution}{exo:b2:integration:8}
$H$ 关于 $x$ 可微（被积函数属于 $C^1$ 类，对 $x$ 而言，导数
$-2x(1+t^2)\cdot\frac{\eu^{-x^2(1+t^2)}}{1+t^2} =
-2x\,\eu^{-x^2}\eu^{-x^2t^2}$，在 $x$ 的紧集上\omterm{def:b2:metric:continuity}{连续}且有界，
对 $t \in \intcc{0}{1}$ 的控制显然成立）：
\[
H'(x) = -2x\,\eu^{-x^2}\int_0^1 \eu^{-x^2t^2}\,\dd t
\overset{u = xt}{=} -2\,\eu^{-x^2}\int_0^x \eu^{-u^2}\,\dd u
= -G'(x),
\]
因为 $G'(x) = 2\eu^{-x^2}\int_0^x \eu^{-t^2}\dd t$（链式法则作用于平方项，
并用微积分基本定理）。于是 $G + H$ 为
常数，等于 $G(0) + H(0) = 0 + \int_0^1 \frac{\dd t}{1+t^2}
= \frac\pi4$。

当 $x \to \infty$ 时：$0 \leq H(x) \leq \eu^{-x^2}\int_0^1 \dd t \to
0$，于是 $G(x) \to \frac\pi4$：
\[
\int_0^\infty \eu^{-t^2}\dd t = \sqrt{\frac\pi4} =
\frac{\sqrt\pi}{2} .
\]
（由此可得 $\Gamma\bigl(\frac12\bigr) = 2\int_0^\infty
\eu^{-t^2}\dd t = \sqrt\pi$，作代换 $t = \sqrt u$。）
\end{solution}

\begin{solution}{exo:b2:integration:9}
柯西--施瓦茨不等式（见本科一年级分册，在 $\intcc{\varepsilon}{M}$
上成立并取极限）应用于分解式
$t^{\frac{x+y}{2}-1}\eu^{-t} = \bigl(t^{x-1}\eu^{-t}\bigr)^{1/2}
\bigl(t^{y-1}\eu^{-t}\bigr)^{1/2}$：
\[
\Gamma\Bigl(\frac{x+y}{2}\Bigr) \leq
\Gamma(x)^{1/2}\,\Gamma(y)^{1/2}
\quad\Longrightarrow\quad
\ln\Gamma\Bigl(\frac{x+y}{2}\Bigr) \leq \frac{\ln\Gamma(x) +
\ln\Gamma(y)}{2} :
\]
$\ln\Gamma$ 是中点凸的；又因其\omterm{def:b2:metric:continuity}{连续}
（由\cref{thm:b2:integration:gammaprops}），它是凸的
（由\cref{exo:b2:realfun:8}）。（对数凸性使 $\Gamma$
在阶乘的各种插值中唯一确定——即
玻尔–莫勒鲁普定理，本科三年级的一颗明珠。）
\end{solution}

\begin{solution}{exo:b2:integration:10}
\begin{enumerate}
  \item 当 $x \geq a > 0$ 时，被积函数关于 $x$ 的导数为
        $-\sin t\,\eu^{-xt}$，被 $\eu^{-at}$ 控制：
        由\cref{thm:b2:integration:leibnizrule}可得 $F'(x) =
        -\int_0^\infty \eu^{-xt}\sin t\,\dd t$。两次分部积分
        （或借助复指数）：
        \[
        \int_0^\infty \eu^{-xt}\sin t\,\dd t
        = \Im \int_0^\infty \eu^{(-x+\iu)t}\dd t
        = \Im\frac{1}{x - \iu} = \frac{1}{1 + x^2} .
        \]
  \item $\abs{F(x)} \leq \int_0^\infty \eu^{-xt}\dd t = \frac1x \to
        0$。将 $F' = -\frac{1}{1+x^2}$ 自 $x$ 积分
        至 $\infty$：$0 - F(x) = -\bigl(\frac\pi2 - \arctan x\bigr)$，
        于是 $F(x) = \frac\pi2 - \arctan x$。
  \item 令 $x \to 0^+$，利用在本阶段不加证明地承认的\omterm{def:b2:metric:continuity}{连续}性：$F(0^+) =
        \frac\pi2$，且 $F(0) = \int_0^\infty \frac{\sin t}{t}\dd
        t$（半收敛的狄利克雷积分，
        见\cref{ex:b2:integration:sint}）：其值为
        $\frac\pi2$。
\end{enumerate}
\end{solution}

\begin{solution}{exo:b2:integration:11}
收敛性：在 $0$ 附近，被积函数可\omterm{def:b2:metric:continuity}{连续}延拓，延拓值为
$1$（$\sin t \sim t$）；无穷远处被积函数 $\leq t^{-2}$，
\omterm{def:b2:series:def}{绝对收敛}。在 $\intcc{\varepsilon}{M}$ 上分部积分，
取 $u = \sin^2 t$，$v' = t^{-2}$：
\[
\int_\varepsilon^M \frac{\sin^2 t}{t^2}\dd t
= \Bigl[-\frac{\sin^2 t}{t}\Bigr]_\varepsilon^M
+ \int_\varepsilon^M \frac{2\sin t\cos t}{t}\dd t
= \Bigl[-\frac{\sin^2 t}{t}\Bigr]_\varepsilon^M
+ \int_{2\varepsilon}^{2M} \frac{\sin u}{u}\dd u
\]
（最后一个积分中 $u = 2t$）。方括号项在两端都
趋于 $0$（$\sin^2\varepsilon/\varepsilon \leq \varepsilon$；
$\sin^2 M/M \leq 1/M$），最后一个积分趋于
$\int_0^\infty \frac{\sin u}{u}\dd u = \frac\pi2$
（由\cref{exo:b2:integration:10}）。于是
\[
\int_0^\infty \Bigl(\frac{\sin t}{t}\Bigr)^{\!2}\dd t
= \frac\pi2 .
\]
\end{solution}

\begin{solution}{exo:b2:integration:12}
\begin{enumerate}
  \item 分部积分，取 $u = \frac{-1}{2t}$，$v' = -2t\,\eu^{-t^2}$
        （从而 $v = \eu^{-t^2}$）：
        \[
        T(x) = \Bigl[\frac{-\eu^{-t^2}}{2t}\Bigr]_x^\infty
        - \int_x^\infty \frac{\eu^{-t^2}}{2t^2}\dd t
        = \frac{\eu^{-x^2}}{2x}
        - \int_x^\infty \frac{\eu^{-t^2}}{2t^2}\dd t .
        \]
        对新的积分用同样的手法（$u = \frac{-1}{4t^3}$，
        $v' = -2t\,\eu^{-t^2}$）：
        \[
        \int_x^\infty \frac{\eu^{-t^2}}{2t^2}\dd t
        = \frac{\eu^{-x^2}}{4x^3}
        - \frac34\int_x^\infty \frac{\eu^{-t^2}}{t^4}\dd t ,
        \]
        由此得到所预告的恒等式。
  \item 再分部积分一次，对余项作估计：
        \[
        \int_x^\infty \frac{\eu^{-t^2}}{t^4}\dd t
        = \frac{\eu^{-x^2}}{2x^5}
        - \frac52\int_x^\infty\frac{\eu^{-t^2}}{t^6}\dd t
        \leq \frac{\eu^{-x^2}}{2x^5},
        \]
        于是 $0 \leq \frac34\int_x^\infty t^{-4}\eu^{-t^2}\dd t
        \leq \frac{3}{8x^5}\eu^{-x^2}$。在第 1 问的恒等式中舍去（正的）
        余项即得下界；舍去（负的）
        第一次分部积分的第二项，
        即得 $T(x) \leq \frac{\eu^{-x^2}}{2x}$。将
        两侧不等式除以 $\frac{\eu^{-x^2}}{2x}$：比值
        夹在 $1 - \frac{1}{2x^2}$ 与 $1$ 之间，于是 $T(x)
        \sim \frac{\eu^{-x^2}}{2x}$。
  \item 反复分部积分，得到形式级数
        \[
        T(x) \approx \frac{\eu^{-x^2}}{2x}\Bigl(1 -
        \frac{1}{2x^2} + \frac{1\cdot3}{(2x^2)^2} -
        \frac{1\cdot3\cdot5}{(2x^2)^3} + \cdots\Bigr),
        \]
        其第 $k$ 个系数 $1\cdot3\cdots(2k-1) =
        \frac{(2k)!}{2^k k!}$ 增长得比任何等比
        数列都快：对固定的 $x$，各项
        $\frac{1\cdot3\cdots(2k-1)}{(2x^2)^k}$ 趋于无穷大
        （其比值为 $\frac{2k+1}{2x^2} \to \infty$），因此
        级数对每个 $x$ 都发散。它是一个
        \emph{渐近}展开式：在任一固定阶截断，
        误差在 $x \to \infty$ 时与第一个被舍去的项同阶，
        ——但它绝非收敛的
        级数。（这一尾项估计即概率各章中的标准高斯
        尾部界。）
\end{enumerate}
\end{solution}

\begin{solution}{pb:b2:integration:1}
\textbf{1.} 在 $0^+$ 处被积函数 $\sim t^{x-1}$，有限端点处的
比较积分收敛当且仅当 $1 - x < 1$，即 $x > 0$
（而 $x \leq 0$ 时 $t^{x-1} \geq t^{-1}$，发散）；在
$+\infty$ 处，$t^{x-1}\eu^{-t} = o(t^{-2})$ 对一切
$x$ 都收敛。于是 $\Gamma(x) = \frac{\Gamma(x+1)}{x}$，且 $\Gamma(x+1)
\to \Gamma(1) = 1$，当 $x \to 0^+$ 时（由\omterm{def:b2:metric:continuity}{连续}性，
\cref{thm:b2:integration:gammaprops}）：$\Gamma(x) \sim
\frac1x$。

\textbf{2.} 令 $t = u^2$，$\dd t = 2u\,\dd u$：
\[
\Gamma\Bigl(\frac12\Bigr) = \int_0^\infty
t^{-1/2}\eu^{-t}\dd t
= \int_0^\infty \frac{\eu^{-u^2}}{u}\,2u\,\dd u
= 2\int_0^\infty \eu^{-u^2}\dd u = \sqrt\pi
\]
由\cref{exo:b2:integration:8}，令 $u = v/\sqrt2$：
\[
\int_\R \eu^{-v^2/2}\dd v
= 2\sqrt2\int_0^\infty \eu^{-u^2}\dd u
= \sqrt2\,\sqrt\pi = \sqrt{2\pi} .
\]

\textbf{3.} $n = 0$ 时成立（两边均为 $\sqrt\pi$）。若
$\Gamma(n + \frac12) = \frac{(2n)!}{4^n n!}\sqrt\pi$，则由
函数方程得
\[
\Gamma\Bigl(n + 1 + \frac12\Bigr)
= \Bigl(n + \frac12\Bigr)\Gamma\Bigl(n + \frac12\Bigr)
= \frac{2n+1}{2}\cdot\frac{(2n)!}{4^n n!}\sqrt\pi
= \frac{(2n+2)!}{4^{n+1}(n+1)!}\sqrt\pi ,
\]
最后一步利用 $\frac{(2n+2)!}{(2n)!} = (2n+2)(2n+1)$
和 $\frac{2n+1}{2} = \frac{(2n+2)(2n+1)}{4(n+1)}$。

\textbf{4.} 由\cref{thm:b2:integration:gammaprops}可得
$\Gamma''(x) = \int_0^\infty t^{x-1}\eu^{-t}(\ln t)^2\dd t$
（两次应用莱布尼茨法则，控制方式与
定理证明中相同）；被积函数 $\geq 0$ 且不恒为
零，故 $\Gamma'' > 0$：$\Gamma$ 严格凸，
$\Gamma'$ 严格递增。由 $\Gamma(1) = \Gamma(2)
= 1$，由罗尔定理得 $x_0 \in \intoo12$，使 $\Gamma'(x_0) =
0$；由 $\Gamma'$ 严格单调，$x_0$ 为其唯一
零点，且 $\Gamma' < 0$ 于其前、$\Gamma' > 0$ 于其后：
$\Gamma$ 在 $\intoo0{x_0}$ 上递减，在
$\intoo{x_0}\infty$ 上递增，且 $x_0$ 是唯一的最小值点。

\textbf{5.} 设 $k \in \N$，$x \geq 3$；取整数
$n$ 满足 $n + 1 \leq x < n + 2$（从而 $n \geq 1$）。由
第 4 题的单调性（自 $x_0 < 2$ 起成立）：$\Gamma(x)
\geq \Gamma(n + 1) = n!$，而 $x^k \leq (n+2)^k$。由于
$\frac{n!}{(n+2)^k} \to \infty$（阶乘增长快于幂，参见
本科一年级分册），故 $\frac{\Gamma(x)}{x^k} \geq \frac{n!}{(n+2)^k} \to
\infty$，当 $x \to \infty$ 时，即 $x^k = o(\Gamma(x))$。

\textbf{6.} 在 $0$ 附近被积函数 $\sim t^{x-1}$
（当且仅当 $x > 0$ 时收敛），在 $1$ 附近它为 $\sim (1-t)^{y-1}$
（当且仅当 $y > 0$）；两个比较都是正函数之间的比较，
故 $B(x,y)$ 恰在 $x, y > 0$ 时收敛。作代换
$t \mapsto 1 - t$ 交换两个因子，故 $B(x,y) = B(y,x)$。

\textbf{7.} $B(x,1) = \int_0^1 t^{x-1}\dd t = \frac1x$。在
$\intcc\varepsilon{1-\varepsilon}$ 上分部积分，取 $u = (1-t)^y$，$v =
\frac{t^x}{x}$：
\[
\int t^{x-1}(1-t)^{y}\dd t
= \Bigl[\frac{t^x(1-t)^y}{x}\Bigr]
+ \frac{y}{x}\int t^{x}(1-t)^{y-1}\dd t ;
\]
方括号在 $\varepsilon \to 0$ 时两端均为零（$x >
0$ 在 $0$ 处，$y > 0$ 在 $1$ 处），于是 $B(x, y+1) = \frac
yx\,B(x+1, y)$。

\textbf{8.} 由 $t + (1-t) = 1$：
\[
t^{x-1}(1-t)^{y-1} = t^{x}(1-t)^{y-1} + t^{x-1}(1-t)^{y},
\]
于是 $B(x,y) = B(x+1,y) + B(x,y+1)$。第 7 题给出 $B(x+1,y) =
\frac xy B(x,y+1)$，代入得
\[
B(x,y) = \Bigl(\frac xy + 1\Bigr)B(x,y+1)
= \frac{x+y}{y}\,B(x,y+1),
\]
即 $B(x,y+1) = \frac{y}{x+y}B(x,y)$；另一关系式
由第 6 题的对称性得出。

\textbf{9.} 对 $n$ 归纳（$m$ 固定）：$B(m,1) = \frac1m =
\frac{(m-1)!\,0!}{m!}$，若公式对 $n$ 成立，
\[
B(m, n+1) = \frac{n}{m+n}\,B(m,n)
= \frac{n}{m+n}\cdot\frac{(m-1)!(n-1)!}{(m+n-1)!}
= \frac{(m-1)!\,n!}{(m+n)!} .
\]
改写得 $B(m,n) = \frac{(m-1)!(n-1)!}{(m+n-1)!} =
\bigl[(m+n-1)\binom{m+n-2}{m-1}\bigr]^{-1}$。

\textbf{10.} 等式 $B(x,n) =
\frac{\Gamma(x)\Gamma(n)}{\Gamma(x+n)}$ 两端均等于 $\frac1x$（当 $n =
1$ 时，$\Gamma(1) = 1$，$\Gamma(x+1) = x\Gamma(x)$）。若两者在
$n$ 处相等，由递降关系与
函数方程得：
\[
B(x, n+1) = \frac{n}{x+n}\,B(x,n), \qquad
\frac{\Gamma(x)\Gamma(n+1)}{\Gamma(x+n+1)}
= \frac{n}{x+n}\cdot
\frac{\Gamma(x)\Gamma(n)}{\Gamma(x+n)} :
\]
两数列满足相同的递推关系，且初值相同，
故对所有 $n \in \N^*$ 及 $x > 0$ 二者相等。

\textbf{11.} 令 $t = \sin^2\theta$（$\theta \in
\intoo0{\pi/2}$, $\dd t = 2\sin\theta\cos\theta\,\dd\theta$),
$t^{x-1} = \sin^{2x-2}\theta$，且 $(1-t)^{y-1} =
\cos^{2y-2}\theta$:
\[
B(x,y) = \int_0^{\pi/2}\sin^{2x-2}\theta\,\cos^{2y-2}\theta
\cdot 2\sin\theta\cos\theta\,\dd\theta
= 2\int_0^{\pi/2}\sin^{2x-1}\theta\,\cos^{2y-1}\theta\,
\dd\theta .
\]

\textbf{12.} 取 $y = \frac12$（使余弦因子消失），
并令 $2x - 1 = n$：$B\bigl(\frac{n+1}2, \frac12\bigr) = 2W_n$，即\
$W_n = \frac12 B\bigl(\frac{n+1}2,\frac12\bigr)$。关于第一变量的
降阶关系给出
\[
\frac{W_n}{W_{n-2}}
= \frac{B\bigl(\frac{n-1}2 + 1, \frac12\bigr)}
{B\bigl(\frac{n-1}2, \frac12\bigr)}
= \frac{\frac{n-1}2}{\frac{n-1}2 + \frac12}
= \frac{n-1}{n} :
\]
即沃利斯递推式，这一次不再对正弦作分部积分——
第二部分已一劳永逸地完成了这项工作。

\textbf{13.} $B\bigl(\frac12,\frac12\bigr) = 2W_0 =
2\cdot\frac\pi2 = \pi$，而
$\Gamma\bigl(\frac12\bigr)^2/\Gamma(1) = (\sqrt\pi)^2 = \pi$:
欧拉公式在 $\bigl(\frac12,\frac12\bigr)$ 处成立。

\textbf{14.} 迭代 $W_{2n} = \frac{2n-1}{2n}W_{2n-2}$，
并以 $W_0 = \frac\pi2$ 为起点：
\[
W_{2n} = \frac\pi2\prod_{k=1}^{n}\frac{2k-1}{2k}
= \frac\pi2\cdot\frac{(2n)!}{4^n(n!)^2},
\]
因为 $\prod(2k-1) = \frac{(2n)!}{2^n n!}$，且 $\prod 2k = 2^n
n!$。于是利用第 3 题：
\[
B\Bigl(n+\frac12, \frac12\Bigr) = 2W_{2n}
= \pi\,\frac{(2n)!}{4^n(n!)^2}
= \frac{(2n)!\sqrt\pi}{4^n n!}\cdot\frac{\sqrt\pi}{n!}
= \frac{\Gamma\bigl(n+\frac12\bigr)\Gamma\bigl(\frac12\bigr)}
{\Gamma(n+1)} .
\]
现固定 $x \in \frac12\N^*$。欧拉公式在 $(x,
\frac12)$ 处成立：$x$ 为整数时由第 10 题（利用对称性）得到，
$x = n + \frac12$ 时即上式。欧拉公式的两端
均满足降阶递推 $y \mapsto y + 1$
（左端由第 8 题，右端由函数方程，
同第 10 题）：由归纳法，公式自 $y =
\frac12$ 与 $y = 1$ 出发，推及每个 $y \in \frac12\N^*$。欧拉
公式因此在 $2x, 2y \in \N^*$ 时成立。

\textbf{15.} 令 $u = \frac{t}{1-t}$，即\ $t =
\frac{u}{1+u}$, $1 - t = \frac{1}{1+u}$, $\dd t =
\frac{\dd u}{(1+u)^2}$:
\[
B(x,y) = \int_0^\infty
\Bigl(\frac{u}{1+u}\Bigr)^{x-1}
\Bigl(\frac{1}{1+u}\Bigr)^{y-1}
\frac{\dd u}{(1+u)^2}
= \int_0^\infty \frac{u^{x-1}}{(1+u)^{x+y}}\,\dd u .
\]
当 $x = y = \frac12$ 时，取 $u = v^2$：
\[
\int_0^\infty \frac{u^{-1/2}}{1+u}\dd u
= \int_0^\infty \frac{2\,\dd v}{1+v^2} = \pi
= B\Bigl(\frac12,\frac12\Bigr) . \checkmark
\]

\textbf{16.} 作一次分部积分，对 $1 \leq k \leq n$
及 $s > 0$（取 $u = (1 - t/n)^k$，$v = t^s/s$；边界项
为零）：
\[
\int_0^n \Bigl(1-\frac tn\Bigr)^{\!k} t^{s-1}\dd t
= \frac{k}{ns}\int_0^n \Bigl(1-\frac tn\Bigr)^{\!k-1}
t^{s}\dd t .
\]
从 $k = n$，$s = x$ 出发，迭代 $n$ 次：
\[
\int_0^n \Bigl(1-\frac tn\Bigr)^{\!n} t^{x-1}\dd t
= \frac{n(n-1)\cdots1}{n^n\,x(x+1)\cdots(x+n-1)}
\int_0^n t^{x+n-1}\dd t
= \frac{n!}{n^n}\cdot
\frac{n^{x+n}}{x(x+1)\cdots(x+n)} ,
\]
即 $\dfrac{n!\,n^x}{x(x+1)\cdots(x+n)}$。

\textbf{17.} 由\cref{exo:b2:integration:7}知左端趋于
$\Gamma(x)$（控制收敛定理，控制函数为
$t^{x-1}\eu^{-t}$）；右端为高斯商：
\[
\Gamma(x) = \lim_{n\to\infty}
\frac{n!\,n^x}{x(x+1)\cdots(x+n)} .
\]

\textbf{18.} 对第 16 题中的商
$G_n(x)$ 取对数，并将 $\ln(x+k) = \ln k + \ln(1 + x/k)$ 拆开，对 $k
\geq 1$:
\[
\ln G_n(x) = x\ln n - \ln x - \sum_{k=1}^n
\ln\Bigl(1+\frac xk\Bigr)
= -\ln x + x(\ln n - H_n)
+ \sum_{k=1}^n\Bigl(\frac xk -
\ln\Bigl(1+\frac xk\Bigr)\Bigr).
\]
当 $u \geq 0$ 时，$u - \frac{u^2}2 \leq \ln(1+u) \leq u$，故
一般项落在 $\intcc{0}{x^2/(2k^2)}$ 中，级数
收敛（与 $\sum k^{-2}$ 比较）。由于 $\ln n - H_n
\to -\gamma$（见\cref{ex:b2:comparison:harmonic}），且 $\ln G_n(x)
\to \ln\Gamma(x)$（由第 17 题及 $\ln$ 的\omterm{def:b2:metric:continuity}{连续}性）：
\[
\ln\Gamma(x) = -\ln x - \gamma x
+ \sum_{k=1}^\infty\Bigl(\frac xk -
\ln\Bigl(1+\frac xk\Bigr)\Bigr) .
\]

\textbf{19.} 当 $x = \frac12$ 时，分母为
$\prod_{k=0}^n\bigl(k+\frac12\bigr) =
\frac{(2n+1)!}{2^{2n+1}n!}$（展开各半整数因子），故
\[
G_n\Bigl(\frac12\Bigr)
= \frac{n!\,\sqrt n\;2^{2n+1}n!}{(2n+1)!}
= \frac{2\sqrt n\;4^n}{(2n+1)\binom{2n}{n}} .
\]
由 $\binom{2n}{n} \sim \frac{4^n}{\sqrt{\pi n}}$
（见\cref{ex:b2:comparison:centralbinomial}），可得：
\[
G_n\Bigl(\frac12\Bigr)
\sim \frac{2\sqrt n\,\sqrt{\pi n}}{2n+1}
\longrightarrow \sqrt\pi = \Gamma\Bigl(\frac12\Bigr) .
\]
中心二项式系数中的 $\sqrt\pi$（来自
沃利斯公式，进而来自斯特林常数），与高斯
积分的 $\sqrt\pi$ 是同一个数。

\textbf{20.} 该恒等式是直接的代数运算：把
$\frac{n!\,n^x}{x(x+1)\cdots(x+n)}$ 乘以 $\frac{nx}{x+n+1}$，并
将 $x$ 并入乘积，$n$ 并入 $n^x$。令 $n \to
\infty$ 时，左边趋于 $\Gamma(x+1)$（高斯公式在 $x+1$ 处），
右边趋于 $\Gamma(x)\cdot x\cdot 1$，因为
$\frac{n}{x+n+1} \to 1$，于是 $\Gamma(x+1) = x\Gamma(x)$——
无需一次分部积分即可推得。

\textbf{21.} 令 $u = t^a$，$t = u^{1/a}$，$\dd t = \frac1a
u^{1/a - 1}\dd u$，于是：
\[
\int_0^\infty \eu^{-t^a}\dd t
= \frac1a\int_0^\infty u^{\frac1a - 1}\eu^{-u}\dd u
= \frac1a\,\Gamma\Bigl(\frac1a\Bigr)
= \Gamma\Bigl(1 + \frac1a\Bigr) .
\]
当 $a \to +\infty$（沿任意数列）时，$\eu^{-t^a} \to 1$ 在
$0 < t < 1$ 上，$\to \eu^{-1}$ 在 $t = 1$ 处，$\to 0$ 在 $t > 1$ 时；
当 $a \geq 2$ 时，取控制函数 $\mathbf 1_{t \leq 1} +
\eu^{-t^2}\mathbf 1_{t > 1}$（$t^a \geq t^2$，当 $t \geq 1$ 时），
它可积。由控制收敛定理，积分趋于
$\int_0^1 1\,\dd t = 1$——，这也必然如此，因为 $\Gamma(1 +
\frac1a) \to \Gamma(1) = 1$，由\omterm{def:b2:metric:continuity}{连续}性得出。

\textbf{22.} 令 $u = t^n$，$\dd t = \frac1n u^{1/n - 1}\dd
u$，于是：
\[
\int_0^1 \frac{\dd t}{\sqrt{1-t^n}}
= \frac1n\int_0^1 u^{\frac1n-1}(1-u)^{-1/2}\dd u
= \frac1n\,B\Bigl(\frac1n, \frac12\Bigr) .
\]
当 $n = 1$ 时，$B\bigl(1,\frac12\bigr) = B\bigl(\frac12,1\bigr) = 2$，
与 $\int_0^1\frac{\dd t}{\sqrt{1-t}} = 2$ 相符。当 $n = 2$ 时，
$\frac12 B\bigl(\frac12,\frac12\bigr) = \frac\pi2 = \arcsin 1$。
当 $n = 4$ 时，$\frac14 B\bigl(\frac14,\frac12\bigr)$ 即
双纽线常数，它不存在初等闭形式。

\textbf{23.} 反复应用函数方程：
\[
\frac{1}{\Gamma(x)}\int_0^\infty t^{x+k-1}\eu^{-t}\dd t
= \frac{\Gamma(x+k)}{\Gamma(x)}
= (x+k-1)(x+k-2)\cdots x ,
\]
这是含 $k$ 个因子的上升阶乘。当 $x = 1$ 时，
$\Gamma(1+k)/\Gamma(1) = k!$，即 $\eu^{-t}$ 的矩，
见\cref{exo:b2:integration:2}。

\textbf{24.} 作代换 $t = \frac{1+s}2$（$s \in
\intoo{-1}1$，$\dd t = \frac{\dd s}2$，$t(1-t) =
\frac{1-s^2}4$）：
\[
B(x,x) = \int_{-1}^{1}\Bigl(\frac{1-s^2}{4}\Bigr)^{x-1}
\frac{\dd s}{2}
= 4^{1-x}\int_0^1 (1-s^2)^{x-1}\dd s
\]
（被积函数为偶函数）。再令 $s = \sqrt v$（$\dd s =
\frac{\dd v}{2\sqrt v}$）：
\[
B(x,x) = \frac{4^{1-x}}{2}\int_0^1
v^{-1/2}(1-v)^{x-1}\dd v
= 2^{1-2x}\,B\Bigl(\frac12, x\Bigr) .
\]
当 $2x \in \N^*$ 时，此处出现的所有参数都落在
$\frac12\N^*$ 中，故欧拉公式（第 14 题）
对两边均适用：
\[
\frac{\Gamma(x)^2}{\Gamma(2x)}
= 2^{1-2x}\,
\frac{\Gamma\bigl(\frac12\bigr)\Gamma(x)}
{\Gamma\bigl(x+\frac12\bigr)}
\quad\Longleftrightarrow\quad
\Gamma(x)\,\Gamma\Bigl(x+\frac12\Bigr)
= 2^{1-2x}\sqrt\pi\;\Gamma(2x) .
\]
当 $x = n$ 时直接验证：左边为 $(n-1)!\cdot
\frac{(2n)!\sqrt\pi}{4^n n!} = \frac{(2n)!\sqrt\pi}{4^n n}$，
右边为 $2\cdot4^{-n}\sqrt\pi\,(2n-1)! =
\frac{(2n)!\sqrt\pi}{4^n n}$，两者相等。

\textbf{25.} (i) 分部积分得到 $B(x,y+1) =
\frac yx B(x+1,y)$，由这一恒等式可推出
所有降阶关系、整数与半整数取值，以及
沃利斯递推式。(ii) 控制收敛定理将
初等积分 $\int_0^n(1-t/n)^n t^{x-1}$ 化为
$\Gamma(x)$（高斯公式，第 17 题），并计算了
第 21 题中 $a \to \infty$ 时的极限。(iii) 由比较一章引入
欧拉常数（$\ln n - H_n \to
-\gamma$，第 18 题），以及中心二项式的渐近
式（第 19 题）——即伪装形式的斯特林公式。
(iv) 欧拉公式 $B(x,y) = \Gamma(x)\Gamma(y)/\Gamma(x+y)$ 现已对
$y \in \N^*$、任意 $x > 0$ 的情形得证（第 10 题），并对全部
半整数对得证（第 14 题）；一般情形
$x, y > 0$ 有待多重积分一章中的二重积分计算，
该计算在四分之一平面上对
$\Gamma(x)\Gamma(y)$ 作乘积分解。
\end{solution}
